% BHU PYQ -- Physics: Applied Electronics, Fiber Optics & Semiconductor Devices (2025-26 NEP)
% Paper Code: PHYMN31 / PHYMJ31 / PHYMV31 (Sem III Re-Mid Term / Sessional Exam, Batch P2)
\documentclass[11pt,a4paper]{article}
\usepackage[a4paper,top=2.5cm,bottom=2.5cm,left=2.8cm,right=2.8cm,headheight=14pt]{geometry}
\usepackage{amsmath,amssymb}
\usepackage[T1]{fontenc}
\usepackage[utf8]{inputenc}
\usepackage{lmodern,microtype,fancyhdr,enumitem,hyperref,lastpage,parskip}

\hypersetup{colorlinks=true,urlcolor=black,linkcolor=black,pdftitle={PHYMN31: Applied Electronics \& Fiber Optics -- BHU Sem III Mid-Term 2025-26 (NEP)}}

\pagestyle{fancy}\fancyhf{}
\renewcommand{\headrulewidth}{0.4pt}\renewcommand{\footrulewidth}{0.4pt}
\fancyhead[L]{\small \textit{Banaras Hindu University}}
\fancyhead[R]{\small \textit{PHYMN31: Applied Electronics (Mid-Term 2025-26)}}
\fancyfoot[C]{\small Page \thepage\ of \pageref{LastPage}}

\newcommand{\pts}[1]{\hfill \textit{[#1]}}

\begin{document}

\begin{center}
  {\large\bfseries BANARAS HINDU UNIVERSITY}\\[2pt]
  {\normalsize B.Sc. (Hons.) Semester III Sessional / Re-Mid Term Examination 2025-26 (NEP)}\\[6pt]
  \rule{0.8\linewidth}{0.5pt}\\[6pt]
  {\large\bfseries DEPARTMENT OF PHYSICS}\\[4pt]
  {\large\bfseries Paper Code: PHYMN31 / PHYMJ31 / PHYMV31 : Applied Electronics and Fiber Optics}\\[2pt]
  {\normalsize (Re-Mid Term Examination --- Batch P2)}\\[4pt]
  {\normalsize Time: 1 Hour \hfill Maximum Marks: 20}
\end{center}
\rule{\linewidth}{0.4pt}
\smallskip

\noindent \textit{Instructions: Answer all questions as specified. Figures to the right indicate full marks.}

\bigskip

\begin{enumerate}[label=\textbf{\arabic*.},leftmargin=*,itemsep=14pt]

  \item How does the refractive index difference between the core and cladding facilitate light transmission in optical fibers?\pts{3}\\
  \begin{center}\textbf{OR}\end{center}
  Why is single-mode fiber preferred for long-distance communication, while multimode fiber is used for shorter connections?\pts{3}

  \item If an optical fiber has a power loss of $5\text{ dB}$, what percentage of the input power is lost?\pts{3}

  \item What are nonlinear losses in optical fibers and how do they differ from linear losses?\pts{4}

  \item A galvanometer has a resistance of $50\ \Omega$ and it deflects full scale when a current of $10\text{ mA}$ flows in it. How can it be converted into an ammeter of range $10\text{ A}$?\pts{3}

  \item What do you mean by intrinsic and extrinsic semiconductor?\pts{2}

  \item An AC voltage source $V(t) = 200 \sin(100 t)\text{ V}$ is connected to a half-wave rectifier with a load resistance of $100\ \Omega$ (resistance of diode is $1\ \Omega$). Calculate:\pts{5}
  \begin{enumerate}[label=(\alph*),itemsep=3pt,topsep=4pt]
    \item DC Power output,
    \item AC input power,
    \item Rectification efficiency
  \end{enumerate}

\end{enumerate}

\bigskip
\begin{center}
  \textbf{***}
\end{center}

\end{document}
